\section{Intuition, Infinity, Science}

Robin Collingwood is probably the last of the idealist philosophers\footnote{\url{https://gornahoor.net/?p=54}} and he was greatly influenced by the Italian philosopher, Giovanni Gentile. Unfortunately, in our time, Gentile is quite forgotten, even by those on the ``right'' who should be recovering his thought. Perhaps these quotes will motivate some to return to the richness of Western thought, rather than worrying about what some Cro-Magnon man was thinking 20,000 years ago. Note, too, the emphasis on will and self-creation. Following the long quote, I offer a brief commentary.

\begin{quotationx}
Man's life is a becoming: and not only becoming, but self-creation. He does not grow under the direction and control of irresistible forces. The force that shapes him is his own will. All his life is an effort to attain to real human nature. But human nature, since man is at bottom spirit, is only exemplified in the absolute spirit of God. Hence man must shape himself in God's image, or he ceases to be even human and becomes diabolical. This self-creation must also be self-knowledge; not the self-knowledge of introspection, the examination of the self that is, but the knowledge of God, the self that is to be. Knowledge of God is the beginning, the center, and end, of human life.

How is the mind to be at once in change and out of change? Only if the mind originates change in itself. For then, as the source and ground of change, it will not be subject to change; while on the other hand, as undergoing change through its own free act, it will exhibit change. This double aspect of the mind as active and passive is the very heart of philosophy the distinction between fact and act . The act is out of time in the sense that it creates time, just as it is supernatural in the sense that it creates nature; the fact is temporal, natural, subject to all those laws which constitute its finiteness. But between the act and the fact there is no division: the distinction is only an ideal distinction. In creating the fact, the act realized itself, and does not live apart in a heaven of its own for which it issues mandates for the creation of facts; it lives in the facts which it creates.

The life of absolute knowledge is thus the conscious self-creation of the mind, no mere discovery of what it is, but the making of itself what it is.

The infinite is not another thing which is best grasped by sweeping the finite out of the way; the infinite is nothing but the unity, or as we sometimes say, the ``meaning,'' of finite things in their diversity and their mutual connections. To look for the infinite by throwing away the finite would be very much like making the players stop playing in order to hear the symphony. What they are collectively playing \emph{is} the symphony; and if you cannot hear it for the noise they are making, you cannot hear it at all. The notes are, so to speak, the body of which the symphony is the soul; and in that sense we might say that the finite is the body of which the infinite is the soul; though, if we say that, we must beware of the materialism which would delude us into talking of disembodies spirits, and remember that it is of the essence of spirit to embody itself.
\flright{\textsc{Robin Collingwood}}
\end{quotationx}


There has been much discussion, and probably more than a little bit of confusion, about the related issues of

\begin{enumerate}


\item Where does science leave off and metaphysics begin?

\item How does the rational knowing of the sciences differ from the intuitional knowing of metaphysics?
\end{enumerate}


I have a short article by the philosopher Robin Collingwood that, I believe, clarifies these issues, even though he frames the discussion as religion (metaphysics) vs. science, and faith (intuition) vs. reason. So, keeping that translation in mind, I can offer a few of his thoughts.

Fundamentally, he points out that faith (intuition or intellect) is the knowledge of the infinite and reason the knowledge of the finite.

What is known by intuition is universal, necessary, and rational, but not demonstrable by logic. It cannot be denied without creating a contradiction. Therefore, it is not knowledge of any particular fact (e.g., my ``intuition'' tells me who is on the telephone). Science deals with the details of the world and a scientific theory is always falsifiable (deniable) in principle. Therefore, it can only be considered \emph{opinion} and lacks the absolute certainty of metaphysical intuition.

The proper sphere of faith (or intellect) is everything in the collective sense --- everything as a whole. The proper sphere of reason is everything in the distributive sense, each separate thing, no matter what. Superstition means to select certain finite things from among the rest and withdraw them from the sphere of reason (that is, it makes the claim that the part is the whole).

Collingwood also deals with the question of why we talk about what allegedly cannot be talked about. If you believe i.e., have faith in, know by intuition in the rationality of the world and the trustworthiness of human thinking, you will embody your belief in detailed scientific inquiries. That shows the foundation of science in an intuition about the whole. Conversely, the spirit of faith intuition is shown to be a real spirit by embodying itself in reason, that is, by developing its own assertions, which as undeveloped would be mere abstractions, into a rational system of thought and conduct.


\flrightit{Posted on 2010-04-05 by Cologero}

\begin{center}* * *\end{center}

\begin{footnotesize}
\begin{sffamily}
\texttt{littleM on 2018-04-05 at 10:50 said:}

I am so glad to be very untalented in fields of empirical science, especially Mathematics. Ontological truths are better supported by simple minds, not speculative conjectures. I am unsure what is the word for it in English, but in German it is Hochmut, in Greek megalopsychia, or Latin superbia -- to dare to ponder and give speculative answers about Prim Truths.

\end{sffamily}
\end{footnotesize}

